\subsection{Complemented Subspaces}

Given a vector space $X$ and a subspace $X_0 \subseteq X$ we may want to decompose $X$ as the direct sum of $X_0$ and some other subspace $X_1 \subseteq X$. More precisely this means that the canonical map 
\[
X_0 \oplus X_1 \to X
\]
induced by the inclusions $X_i \hookrightarrow X$ should be an isomorphism. In the category of vector spaces this is of course always possible and we call $X_1$ and \textbf{algebraic complement} of $X_0$. However the situation looks quite a bit different in the category of topological vector spaces, where such $X_1$ may not always exist. If it exists we call $X_1$ a \textbf{topological complement} of $X_0$ and the space $X_0$ \textbf{complemented}. The following Lemma gives a characterization for when an algebraic complement is also a topological complement (at least in the case of Banach spaces).

\begin{lemma}
    \label{funkana:lem:alg_vs_top_complement}
    Let $X$ be a Banach space and $X_0,X_1 \subseteq X$ subspaces, where $X_1$ is an algebraic complement of $X_0$. Then $X_1$ is a topological complement of $X_0$ if and only if $X_0$ and $X_1$ are closed.
\end{lemma}
\begin{proof}
    Clearly $X_0,X_1 \subseteq X_0 \oplus X_1$ are closed, so if $X$ is toplinearly isomorphic to $X_0 \oplus X_1$ they are also closed in $X$. On the contrary if $X_0$ and $X_1$ are closed, they are again Banach spaces and since we assume the canonical map 
    \[
    X_0 \oplus X_1 \to X
    \]
    to be a bijection it follows from the open mapping theorem that this map is already a toplinear isomorphism.
\end{proof}

Here are two simple sufficient conditions for the existence of topological complements. For a vector space $X$ and a subspace $X_0 \subseteq X$ we define its \textbf{codimension} by 
\[
\mathrm{codim}(X_0) := \mathrm{dim}\left(\modulo{X}{X_0}\right).
\]

\begin{lemma}
    Let $X$ be a Banach space and $X_0 \subseteq X$ a subspace. The following two conditions are both sufficient for the existence of a topological complement of $X_0$: 
    \begin{enumerate}[(i)]
        \item $\mathrm{dim}(X_0) < \infty$
        \item $\mathrm{codim}(X_0) < \infty$ and $X_0$ is closed
    \end{enumerate}
\end{lemma}

\begin{proof}
    First assume $(i)$. Choose a basis $e_1,\ldots,e_n$ of $X_0$ and define linear maps $\phi_i \colon X_0 \to \R$ by 
    \[
    \phi_i(e_j) = \delta_{ij}.
    \]
    Since $X_0$ is finite dimensional the $\phi_i$ are continuous. So using the Hahn-Banach theorem we can extend them to continuous linear functionals $\Phi_i$ defined on all of $X$. Now set 
    \[
    X_1 := \{x \in X \mid \Phi_1(x) = \ldots = \Phi_n(x) = 0\}.
    \]
    Clearly $X_1\subseteq X$ is closed, so by Lemma \ref{funkana:lem:alg_vs_top_complement} we only have to show that $X_1$ is an algebraic complement of $X_0$. To this end let $x \in X$ be arbitary and define 
    \[
    x_0 := \sum_{i=1}^{n} \Phi_i(x)e_i \in X_0.
    \]
    Then $x_1 := x - x_0 \in X_1$ and $x = x_0 + x_1$. Finally observe that $X_0 \cap X_1 = 0$ from which the statement follows. 

    For $(ii)$ we begin by choosing any algebraic complement $X_1$ of $X_0$. Since the restriction of the projection
    \[
    X \to \modulo{X}{X_0}
    \]
    to $X_1$ is an isomorphism, our assumption tells us that $X_1$ is finite dimensional and so in particular closed. Thus the statement follows again from Lemma \ref{funkana:lem:alg_vs_top_complement}.
\end{proof}


